\documentclass[12pt, twoside]{article}
%\documentclass[12pt, twoside]{article}
\usepackage[letterpaper, margin=1in, headsep=0.2in]{geometry}
\setlength{\headheight}{0.6in}
%\usepackage[english]{babel}
\usepackage[utf8]{inputenc}
\usepackage{microtype}
\usepackage{amsmath}
\usepackage{amssymb}
%\usepackage{amsfonts}
\usepackage[nomessages]{fp} %\FPeval{\var-name}{2*sin(pi/6)}
\usepackage{siunitx} %units in math. eg 20\milli\meter
\usepackage{yhmath} % for arcs, overparenth command
\usepackage{tikz} %graphics
\usetikzlibrary{quotes, angles, arrows, arrows.meta}
\usepackage{pgfplots}
\usepgfplotslibrary{fillbetween}
\pgfplotsset{compat=1.16}
\usepackage{graphicx} %consider setting \graphicspath{{images/}}
\usepackage{parskip} %no paragraph indent
\usepackage{enumitem}
\usepackage{multicol}
\usepackage{venndiagram}

\usepackage{fancyhdr}
\pagestyle{fancy}
\fancyhf{}
\renewcommand{\headrulewidth}{0pt} % disable the underline of the header
\raggedbottom
\hfuzz=2mm %suppresses overfull box warnings

\usepackage{hyperref}

\fancyhead[LO]{La Scuola d'Italia / Dr.\ Huson / IB Math 11 \\* 17 November 2026}
\fancyhead[RO]{\fbox{\begin{minipage}{6cm}\footnotesize Nickname:\\[0.5cm]\end{minipage}}}
\fancyfoot[C]{\thepage}

\begin{document}
\subsubsection*{3.2 Pre-Quiz: Sets, Venn Diagrams and Review}

A calculator is allowed. When a question asks for an exact answer or a
particular form, show the algebra that leads to it.

\begin{enumerate}[itemsep=0.15cm]

%--- Problem 1: Three-set Venn word problem: regions in terms of x from set and pairwise totals; solve for x; exactly one; probabilities [13 marks] ---
%--- Source: Original (AA SL prior learning / 4.5 Venn diagrams); three-set extension of 3-1CW_Sets-Venn-review P8 ---
\item[1.]
A group of $70$ students were asked which of three clubs they belong to:
Chess ($C$), Drama ($D$) and Robotics ($R$).

\begin{itemize}[itemsep=0cm, topsep=0.1cm]
  \item $33$ belong to Chess, $24$ belong to Drama and $34$ belong to Robotics.
  \item $11$ belong to both Chess and Drama, $14$ belong to both Chess and
  Robotics, and $8$ belong to both Drama and Robotics.
  \item $x$ students belong to all three clubs, and $7$ students belong to
  none of the three clubs.
\end{itemize}

\begin{enumerate}[itemsep=0.1cm]
  \item Complete the Venn diagram. Write the number of students in each
  region, in terms of $x$ where needed. Start with $x$ in the centre.
\end{enumerate}

\noindent
\begin{minipage}[t]{0.5\textwidth}
\vspace{0pt}
\begin{venndiagram3sets}[labelA=$C$, labelB=$D$, labelC=$R$, tikzoptions={scale=1.3}]
\end{venndiagram3sets}
\end{minipage}%
\begin{minipage}[t]{0.5\textwidth}
\vspace{0pt}
\begin{enumerate}[itemsep=0.1cm, leftmargin=0.9cm]
  \item[(b)] Write down an equation in $x$, and solve it.
  \item[(c)] Find the number of students who belong to exactly one club.
  \item[(d)] A student is chosen at random from the group. Find the
  probability that the student belongs to
  \begin{enumerate}[label=(\roman*), itemsep=0cm]
    \item at least two of the clubs;
    \item Robotics but not Chess.
  \end{enumerate}
\end{enumerate}
\end{minipage}

\medskip

\begin{tikzpicture}
  \draw (-0.6,0) rectangle (14.6,6.0);
  \foreach \y in {1,2,...,5} \draw[dotted] (0,\y)--(14.1,\y);
\end{tikzpicture}

\newpage

%--- Problem 2: Two-set percentages: both from "neither"; probability of cat-not-dog, exactly one, not a dog; expected count; reverse problem P(cat) from dog, both, neither [12 marks] ---
%--- Source: Original; parallel of 3-1CW_Sets-Venn-review P9 (modelled on AA SL 2021 May P2 TZ2 Q4) ---
\item[2.]
In a town, $58\%$ of the households have a dog ($D$) and $35\%$ have a cat
($C$). $24\%$ of the households have neither a dog nor a cat.

\begin{enumerate}[itemsep=0.1cm]
  \item Draw a Venn diagram, and find the percentage of households that have
  both a dog and a cat.
  \item A household is chosen at random. Find the probability that the
  household
  \begin{enumerate}[label=(\roman*), itemsep=0cm]
    \item has a cat but not a dog;
    \item has exactly one of the two pets;
    \item does not have a dog.
  \end{enumerate}
  \item There are $1200$ households in the town. Find the number of
  households that have both a dog and a cat.
  \item In a second town, $46\%$ of the households have a dog, $15\%$ have
  both a dog and a cat, and $32\%$ have neither. A household in this town
  is chosen at random. Find the probability that it has a cat.
\end{enumerate}

\begin{tikzpicture}
  \draw (-0.6,0) rectangle (14.6,12.4);
  \foreach \y in {1,2,...,11} \draw[dotted] (0,\y)--(14.1,\y);
\end{tikzpicture}

\newpage

%--- Problem 3: Smoothie prices: set up a 3x3 system from two sales totals and a price difference; solve; number of medium smoothies [8 marks] ---
%--- Source: Original (AA SL 1.1 / prior learning, systems of linear equations); 3x3 extension of 2-5PQ_Sequences-lines-exponents P4 ---
\item[3.]
A caf\'e sells small, medium and large smoothies. A small smoothie costs
$x$ euros, a medium costs $y$ euros and a large costs $z$ euros.

\begin{itemize}[itemsep=0cm, topsep=0.1cm]
  \item On Monday the caf\'e sold $12$ small, $9$ medium and $6$ large
  smoothies, for a total of $107.10$ euros.
  \item On Tuesday it sold $8$ small, $14$ medium and $5$ large smoothies,
  for a total of $109.10$ euros.
  \item A large smoothie costs $1.50$ euros more than a small smoothie.
\end{itemize}

\begin{enumerate}[itemsep=0.1cm]
  \item Write down three equations in $x$, $y$ and $z$.
  \item Find the price of each size of smoothie.
  \item On Wednesday the caf\'e sold $15$ small and $10$ large smoothies,
  and some medium smoothies, for a total of $153.30$ euros. Find the
  number of medium smoothies sold on Wednesday.
\end{enumerate}

\begin{tikzpicture}
  \draw (-0.6,0) rectangle (14.6,13.2);
  \foreach \y in {1,2,...,12} \draw[dotted] (0,\y)--(14.1,\y);
\end{tikzpicture}

\newpage

%--- Problem 4: Surds and powers: integers a, b from a surd product; hence a quotient; same-base equation; equation quadratic in 3^x [14 marks] ---
%--- Source: Original (AA SL 1.5, 1.7, prior learning surds); in the style of 2-6Q early finishers (a) P1 and 2-10CW P4-P5 ---
\item[4.]
In this question, give every answer in exact form.

\begin{enumerate}[itemsep=0.1cm]
  \item Find the integers $a$ and $b$ such that
  $\left(4 + \sqrt{3}\right)\left(a + b\sqrt{3}\right) = 14 - 3\sqrt{3}$.
  \item Hence write $\dfrac{14 - 3\sqrt{3}}{4 + \sqrt{3}}$ in the form
  $p + q\sqrt{3}$, where $p, q \in \mathbb{Z}$.
  \item Solve the equation $8^{x + 1} = 16^{2x - 1}$.
  \item By letting $u = 3^{x}$, solve the equation
  $3^{2x + 1} - 28\left(3^{x}\right) + 9 = 0$.
\end{enumerate}

\begin{tikzpicture}
  \draw (-0.6,0) rectangle (14.6,16.5);
  \foreach \y in {1,2,...,15} \draw[dotted] (0,\y)--(14.1,\y);
\end{tikzpicture}

\end{enumerate}
\end{document}
